\chapter{Sampling Distributions, Confidence Intervals, and Bootstrap}
\chaptermark{Confidence Intervals and Bootstrap}
\label{chap:confidence-intervals}

\begingroup
\clubpenalty=10000
\widowpenalty=10000
% Shared Chapter 6 scientific charts: analytic curves and reproducible
% simulations. Regenerate the committed tables with generate_chapter6_figures.py.
% The conspect and slides may override these dimensions and the data path.
\providecommand{\InfDataPath}{figures/data}
\providecommand{\InfPlotWidth}{11.8cm}
\providecommand{\InfPlotHeight}{4.6cm}
\definecolor{InfBlue}{HTML}{1F4E79}
\definecolor{InfOrange}{HTML}{C46A13}
\definecolor{InfTeal}{HTML}{007D80}
\definecolor{InfMuted}{HTML}{526477}
\pgfplotsset{inf axis/.style={
  scale only axis,width=\InfPlotWidth,height=\InfPlotHeight,
  axis lines=left,enlargelimits=false,
  grid=major,grid style={black!9},
  tick label style={font=\footnotesize},label style={font=\footnotesize},
  legend style={draw=none,fill=none,font=\footnotesize,
    at={(0.5,1.05)},anchor=south,legend columns=-1,column sep=0.8em},
  every axis plot/.append style={no marks,line width=1.15pt},
}}

\newcommand{\InfCoveragePlot}{%
\begin{tikzpicture}\begin{axis}[inf axis,
  xmin=7.8,xmax=12.4,ymin=0,ymax=41,
  xlabel={Candidate population mean $\mu$},ylabel={Repeated dataset},
  xtick={8,9,10,11,12},ytick={1,10,20,30,40},
  ymajorgrids=false]
  \addplot[InfMuted,densely dashed,forget plot]
    coordinates {(10,0) (10,41)};
  \addplot[InfTeal,only marks,mark=*,mark size=0.65pt,
    error bars/.cd,x dir=both,x explicit,
    error bar style={line width=0.6pt},error mark options={rotate=90,mark size=1.0pt}]
    table[x=mean,y=trial,x error=halfwidth] {\InfDataPath/ch6-coverage-hit.dat};
  \addlegendentry{Contains $\mu=10$}
  \addplot[InfOrange,only marks,mark=*,mark size=0.85pt,
    error bars/.cd,x dir=both,x explicit,
    error bar style={line width=1.0pt},error mark options={rotate=90,mark size=1.2pt}]
    table[x=mean,y=trial,x error=halfwidth] {\InfDataPath/ch6-coverage-miss.dat};
  \addlegendentry{Misses $\mu=10$}
\end{axis}\end{tikzpicture}%
}

\newcommand{\InfStudentPlot}{%
\begin{tikzpicture}\begin{axis}[inf axis,
  xmin=-4,xmax=4,ymin=0,ymax=0.425,
  xlabel={Standardized estimation error},ylabel={Density},
  xtick={-4,-2,0,2,4},ytick={0,0.1,0.2,0.3,0.4}]
  \addplot[draw=none,fill=InfOrange!18,forget plot]
    table[x=x,y=lefttail] {\InfDataPath/ch6-student.dat} \closedcycle;
  \addplot[draw=none,fill=InfOrange!18,forget plot]
    table[x=x,y=righttail] {\InfDataPath/ch6-student.dat} \closedcycle;
  \addplot[InfBlue] table[x=x,y=normal] {\InfDataPath/ch6-student.dat};
  \addlegendentry{Standard Normal}
  \addplot[InfOrange,dashed] table[x=x,y=t9] {\InfDataPath/ch6-student.dat};
  \addlegendentry{Student $t_9$}
  \addplot[InfBlue,densely dotted,forget plot]
    coordinates {(-1.959963985,0) (-1.959963985,0.05844507)};
  \addplot[InfBlue,densely dotted,forget plot]
    coordinates {(1.959963985,0) (1.959963985,0.05844507)};
  \addplot[InfOrange,densely dotted,forget plot]
    coordinates {(-2.262157163,0) (-2.262157163,0.0407427)};
  \addplot[InfOrange,densely dotted,forget plot]
    coordinates {(2.262157163,0) (2.262157163,0.0407427)};
  \node[anchor=north east,font=\footnotesize,text=InfMuted,align=right,
    fill=white,inner sep=2pt] at (axis description cs:0.98,0.97)
    {$z_{0.975}=1.960$\\$t_{9,0.975}=2.262$};
\end{axis}\end{tikzpicture}%
}

\newcommand{\InfProportionCoveragePlot}{%
\begin{tikzpicture}\begin{axis}[inf axis,
  xmin=0,xmax=1,ymin=0,ymax=1.03,
  xlabel={True Bernoulli probability $p$},ylabel={Exact coverage probability},
  xtick={0,0.2,0.4,0.6,0.8,1},ytick={0,0.25,0.5,0.75,0.95}]
  \addplot[InfMuted,densely dotted]
    coordinates {(0,0.95) (1,0.95)};
  \addlegendentry{Nominal $95\%$}
  \addplot[InfOrange,line width=0.85pt]
    table[x=p,y=wald] {\InfDataPath/ch6-proportion-coverage.dat};
  \addlegendentry{Wald, $n=20$}
  \addplot[InfTeal,line width=1.0pt]
    table[x=p,y=wilson] {\InfDataPath/ch6-proportion-coverage.dat};
  \addlegendentry{Wilson, $n=20$}
\end{axis}\end{tikzpicture}%
}

\newcommand{\InfBootstrapPlot}{%
\begin{tikzpicture}\begin{axis}[inf axis,
  xmin=2.3,xmax=7,ymin=0,ymax=0.75,
  xlabel={Bootstrap sample mean $\bar X^*$},ylabel={Bootstrap density},
  xtick={2.5,3.4,4.5,5.68,7},
  xticklabels={$2.5$,$3.40$,$4.5$,$5.68$,$7$}]
  \addplot[InfBlue,fill=InfBlue!20,const plot,line width=0.65pt]
    table[x=value,y=density] {\InfDataPath/ch6-bootstrap-mean.dat} \closedcycle;
  \addlegendentry{$20\,000$ resamples of $n=40$}
  \addplot[InfOrange,dashed,forget plot]
    coordinates {(3.399081569,0) (3.399081569,0.7)};
  \addplot[InfOrange,dashed,forget plot]
    coordinates {(5.681750915,0) (5.681750915,0.7)};
  \addplot[InfTeal,densely dotted,forget plot]
    coordinates {(4.494628920,0) (4.494628920,0.7)};
  \node[anchor=north east,font=\footnotesize,text=InfMuted,align=right,
    fill=white,inner sep=2pt] at (axis description cs:0.98,0.98)
    {Observed mean: $4.495$\\Percentile CI: $[3.40,5.68]$};
\end{axis}\end{tikzpicture}%
}

\newcommand{\InfPairedPlot}{%
\begin{tikzpicture}\begin{axis}[inf axis,
  xmin=-0.02,xmax=0.23,ymin=0,ymax=0.072,
  xlabel={Bootstrap accuracy difference $\widehat a_A^*-\widehat a_B^*$},
  ylabel={Bootstrap probability},xtick={0,0.04,0.1,0.16,0.2},
  xticklabels={$0$,$0.04$,$0.10$,$0.16$,$0.20$},
  ytick={0,0.02,0.04,0.06},scaled y ticks=false,
  yticklabel={\pgfmathprintnumber[fixed,precision=2]{\tick}}]
  \addplot[InfBlue,ycomb,line width=1.7pt]
    table[x=difference,y=mass] {\InfDataPath/ch6-paired-accuracy.dat};
  \addlegendentry{Paired bootstrap, $n=200$, $B=20\,000$}
  \addplot[InfMuted,densely dotted,forget plot]
    coordinates {(0,0) (0,0.068)};
  \addplot[InfOrange,dashed,forget plot]
    coordinates {(0.04,0) (0.04,0.068)};
  \addplot[InfOrange,dashed,forget plot]
    coordinates {(0.16,0) (0.16,0.068)};
  \node[anchor=north east,font=\footnotesize,text=InfMuted,align=right,
    fill=white,inner sep=2pt] at (axis description cs:0.98,0.98)
    {Observed gap: $0.10$\\Percentile CI: $[0.04,0.16]$};
\end{axis}\end{tikzpicture}%
}


\section{What uncertainty remains after fitting}
\label{sec:interval-uncertainty}

The five waiting times from Chapter~\ref{chap:estimation},
\[
  1,\ 2,\ 3,\ 4,\ 5\ \text{minutes},
\]
gave an exponential rate estimate of \(\widehat\lambda=1/3\) requests
per minute. This is an estimate from just five observations. A second
collection of five waits would usually give a different rate. How much
variation should we expect, and how can we report it alongside the estimate?

Section~\ref{sec:samples} already introduced sampling distributions and
standard errors. We now apply that reasoning to the estimators from
Chapter~\ref{chap:estimation}. To study a rule
\(\widehat\theta=T(X_1,\ldots,X_n)\), imagine keeping the population
parameter \(\theta\) fixed, drawing a new dataset, and applying the
\emph{same rule} each time. A histogram of the resulting estimates describes
how the estimation procedure varies across datasets.

Three distributions answer different questions. The observation law
describes a new wait or measurement. A sampling distribution describes
the estimate obtained from a new dataset. A posterior, as in
Section~\ref{sec:map-estimation}, describes uncertainty about a parameter
conditional on the recorded data and a specified prior. The calculations
below use sampling distributions to build intervals with a stated
frequency of covering the true parameter.

For a population mean, the earlier formulas give
\[
  \operatorname{SE}(\bar X_n)=\frac{\sigma}{\sqrt n},
  \qquad
  \widehat{\operatorname{SE}}(\bar X_n)=\frac{S_n}{\sqrt n}.
\]
The second expression estimates the first when \(\sigma\) is unknown.
To turn an error scale into an interval, we also need to know, exactly
or approximately, how often errors exceed a multiple of that scale.

\section{Confidence intervals and repeated sampling}
\label{sec:confidence-coverage}

Suppose a sensor measures an unknown constant level \(\mu\). Each
dataset produces a mean and an interval around it. Across repeated
experiments, both endpoints move. A useful interval rule should contain
\(\mu\) in a specified proportion of those experiments.

\begin{definition}[Confidence interval and coverage]
Let \(0<\alpha<1\). A confidence interval procedure reports
\(C(X)=[L(X),U(X)]\), where both endpoints are functions of the sample.
Its coverage at parameter \(\theta\) is
\[
  \Prob_\theta\{L(X)\le\theta\le U(X)\}.
\]
It has confidence level at least \(1-\alpha\) if this probability is
at least \(1-\alpha\) for every parameter allowed by the model.
An exact central construction may have equality. An approximate
\(95\%\) interval targets coverage \(0.95\) using an approximation.
\end{definition}

The probability concerns the random sample and the interval it produces.
After we observe a particular dataset, its numerical endpoints are fixed.
The frequentist model also treats \(\theta\) as fixed, so its coverage
statement does not supply a posterior probability for this one interval.
For example, saying that a procedure covers the true sensor level in
\(95\%\) of repeated experiments explains the guarantee that we used
when constructing the reported interval.

\begin{figure}[H]
  \centering\InfCoveragePlot
  \caption{Repeated \(95\%\) Gaussian mean intervals. Each horizontal
  segment comes from a different dataset, with \(n=25\), true mean
  \(\mu=10\), and known observation standard deviation \(\sigma=2\).
  The dashed line marks the same true mean in every experiment. Orange
  intervals miss it. A finite collection need not contain exactly
  \(95\%\) successful intervals.}
  \label{fig:interval-coverage}
\end{figure}

\subsection{A Gaussian mean with known variance}

For independent Gaussian observations with mean \(\mu\) and known
variance \(\sigma^2>0\), the exact Gaussian sampling law from
Section~\ref{sec:clt} gives
\[
  Z=\frac{\bar X_n-\mu}{\sigma/\sqrt n}\sim\mathcal N(0,1).
\]
The numerator is the estimation error. Dividing by its standard error
puts it on a scale whose distribution is the same for every value of
the unknown \(\mu\).

\begin{definition}[Quantile and pivot]
For a random variable with CDF \(F\), its \(p\)-quantile is
\(q_p=\inf\{x:F(x)\ge p\}\), for \(0<p<1\).
Write \(z_p\) for the standard Gaussian quantile.
A \emph{pivot} is a function of the observations and the unknown
parameter whose sampling distribution does not depend on unknown
parameters. The variable \(Z\) above is a pivot.
\end{definition}

The standard Gaussian places \(1-\alpha\) of its probability between
\(-z_{1-\alpha/2}\) and \(z_{1-\alpha/2}\). Therefore
\[
  \Prob_\mu\!\left\{
  -z_{1-\alpha/2}\le\frac{\bar X_n-\mu}{\sigma/\sqrt n}
  \le z_{1-\alpha/2}\right\}=1-\alpha.
\]
Rearranging the inequalities leaves the unknown \(\mu\) between
two quantities that we can compute from the data.

\begin{result}[Gaussian mean interval with known variance]
For an i.i.d. Gaussian sample with known \(\sigma^2>0\), the interval
\[
  \left[\bar X_n-z_{1-\alpha/2}\frac{\sigma}{\sqrt n},\quad
        \bar X_n+z_{1-\alpha/2}\frac{\sigma}{\sqrt n}\right]
\]
has exact coverage \(1-\alpha\). For \(95\%\) confidence,
\(z_{0.975}\approx1.96\).
\end{result}

For example, if \(n=25\), \(\bar x=10.4\), and a known
\(\sigma=2\), the standard error is \(0.4\). The interval is
\[
  10.4\pm1.96(0.4)=[9.616,\,11.184].
\]
Its endpoints and \(\mu\) share the measurement units.
The number \(1.96\) is dimensionless.

\subsection{Interval width and the sampling design}

The margin of error in this construction is
\(m=z_{1-\alpha/2}\sigma/\sqrt n\). More observation noise increases
the margin. More independent observations reduce it. Raising the
confidence level increases the quantile and widens the interval.
For example, the Gaussian multipliers are approximately \(1.645\)
at \(90\%\), \(1.960\) at \(95\%\), and \(2.576\) at \(99\%\).

For known \(\sigma\), achieving a chosen margin \(m>0\) requires
\[
  n\ge\left(\frac{z_{1-\alpha/2}\sigma}{m}\right)^2,
\]
rounded upward to an integer. Thus halving the margin requires four
times as many independent observations. This calculation inherits the
sampling assumptions. Correlated measurements do not necessarily
provide that improvement, as Section~\ref{sec:averaging-failures} showed.

\subsection{What the interval refers to}

The interval estimates a population parameter, here the sensor's mean.
It need not contain \(95\%\) of individual readings. Even if we measure
the mean very accurately, a new reading still has observation noise.
Section~\ref{sec:mean-intervals} will make this distinction explicit.

A Bayesian credible interval answers a different probability question:
it contains a specified amount of the posterior probability after
conditioning on the observations and a prior. The posterior calculation
in Section~\ref{sec:map-estimation} supplies that distribution. A
confidence interval derives its guarantee from repetition of the
sampling procedure. The two approaches can sometimes produce the same
numerical interval, but they justify it differently.

Coverage also depends on choosing the procedure as specified. Reporting
only unusually narrow intervals, repeatedly inspecting fresh data until
an interval gives a desired conclusion, or selecting a model after
seeing the same test results changes the experiment. The original
fixed-sample coverage statement alone does not justify such a procedure.

\section{Mean intervals with unknown variance}
\label{sec:mean-intervals}

In most applications \(\sigma\) is unknown. We estimate it by \(S_n\),
using the sample variance from Section~\ref{sec:samples} with denominator
\(n-1\). For small samples, this replacement adds appreciable uncertainty:
some datasets happen to have a small observed spread and hence an
overly optimistic estimated standard error. A wider reference
distribution accounts for that fluctuation under a Gaussian model.

\subsection{Why Student's distribution appears}

We already met Student's \(t\)-distribution as a Gaussian with a random
precision in Section~\ref{sec:student}. Here its random denominator
comes from estimating the observation variance.

\begin{definition}[Chi-square distribution]
If \(Z_1,\ldots,Z_\nu\) are independent standard Gaussians, then
\[
  U=\sum_{j=1}^{\nu}Z_j^2\sim\chi^2_\nu.
\]
The integer \(\nu\) counts the independent squared Gaussian components
and is called the number of degrees of freedom. In the rate convention
of Section~\ref{sec:gamma}, this law is
\(\operatorname{Gamma}(\nu/2,1/2)\).
\end{definition}

For an i.i.d. Gaussian sample, the fitted deviations
\(X_i-\bar X_n\) sum to zero. Estimating one mean leaves \(n-1\)
independent directions of Gaussian variation. More precisely, a Gaussian
sampling result states that
\[
  U=\frac{(n-1)S_n^2}{\sigma^2}\sim\chi^2_{n-1},
  \qquad
  Z=\frac{\sqrt n(\bar X_n-\mu)}{\sigma}\sim\mathcal N(0,1),
\]
and \(U\) and \(Z\) are independent. The independence uses the Gaussian
assumption. It follows by separating the mean direction from the
orthogonal residual directions in a Gaussian vector, using the
zero-covariance independence property from
Section~\ref{sec:multivariate-gaussian}.

Consequently, for \(\nu=n-1\),
\[
  \frac{Z}{\sqrt{U/\nu}}
  =\frac{\sqrt n(\bar X_n-\mu)}{S_n}.
\]
Conditional on \(U\), the left-hand side is Gaussian with precision
\(U/\nu\). This precision has law
\(\operatorname{Gamma}(\nu/2,\nu/2)\), exactly the mixture construction
that previously produced \(t_\nu\).

\begin{theorem}[Student pivot for a Gaussian mean]
For \(n\ge2\) i.i.d. observations from
\(\mathcal N(\mu,\sigma^2)\), with \(\sigma^2>0\),
\[
  \frac{\bar X_n-\mu}{S_n/\sqrt n}\sim t_{n-1}.
\]
If \(t_{\nu,p}\) is the \(p\)-quantile of \(t_\nu\), then
\[
  \left[\bar X_n-t_{n-1,1-\alpha/2}\frac{S_n}{\sqrt n},\quad
        \bar X_n+t_{n-1,1-\alpha/2}\frac{S_n}{\sqrt n}\right]
\]
has exact coverage \(1-\alpha\).
\end{theorem}

The derivation repeats the Gaussian pivot inversion, with a Student
quantile in place of a Gaussian quantile. At low degrees of freedom,
Student's heavier tails require a larger multiplier. As the degrees
of freedom grow, the multiplier approaches the Gaussian one.

\begin{figure}[H]
  \centering\InfStudentPlot
  \caption{A standard Gaussian and Student's \(t_9\) sampling reference
  densities. The Student reference has heavier tails. Its \(97.5\%\)
  quantile is approximately \(2.262\), compared with \(1.960\) for the
  Gaussian. Both intervals allocate \(2.5\%\) to each tail of their
  own reference distribution.}
  \label{fig:interval-student}
\end{figure}

For ten Gaussian measurements with observed mean \(\bar x=10.2\)
and sample standard deviation \(s=1.5\),
\[
  \widehat{\operatorname{SE}}=\frac{1.5}{\sqrt{10}}\approx0.4743,
  \qquad t_{9,0.975}\approx2.262.
\]
The \(95\%\) interval is
\[
  10.2\pm2.262(0.4743)\approx[9.127,\,11.273].
\]
Using \(1.96\) with this estimated standard error would give the
narrower interval \([9.270,11.130]\). Under the Gaussian model at
this sample size, that narrower procedure falls short of \(95\%\)
coverage.

\subsection{A large-sample interval for non-Gaussian observations}

Gaussian sampling gives an exact Student law. For general i.i.d.
observations with finite mean and finite, positive variance, use the
CLT from Section~\ref{sec:clt}. The LLN applied to \(X_i\) and
\(X_i^2\) also makes \(S_n\) approach \(\sigma\). Replacing
\(\sigma\) by this consistent estimate preserves the limiting law:
\[
  \frac{\sqrt n(\bar X_n-\mu)}{S_n}
  \xrightarrow{d}\mathcal N(0,1).
\]
This gives an approximate interval
\(\bar x\pm z_{1-\alpha/2}s/\sqrt n\). Using a Student multiplier
is also common and has the same large-sample limit, but does not make
the interval exact for non-Gaussian data.

The accuracy depends on the observation law and sample size. Strong
skewness and rare extreme values can make convergence slow. There is
no universal sample size at which every population has a reliable
Gaussian approximation. Infinite variance and dependence require
separate analysis rather than automatic use of this formula.

\subsection{A mean interval and a prediction interval}

Under the same i.i.d. Gaussian model, consider a future independent
reading \(X_{\mathrm{new}}\). Its difference from the current sample
mean has variance
\[
  \Var(X_{\mathrm{new}}-\bar X_n)
  =\sigma^2+\frac{\sigma^2}{n}.
\]
The observation noise remains even when we know the mean accurately.
The corresponding exact \(1-\alpha\) prediction interval is
\[
  \bar x\pm t_{n-1,1-\alpha/2}\,s\sqrt{1+\frac1n}.
\]
Its coverage concerns the joint experiment of the training sample
and a new reading. The mean interval uses \(s/\sqrt n\) instead.
As \(n\) grows, the mean interval shrinks, while the prediction
interval retains the scale of individual observation noise.

\subsection{Another exact pivot: an exponential rate}

Return to the five waits at the beginning. The Gamma sum result gives,
for the random total \(W=\sum_iT_i\),
\[
  W\sim\operatorname{Gamma}(n,\lambda),\qquad
  2\lambda W\sim\operatorname{Gamma}(n,1/2)=\chi^2_{2n}.
\]
Let \(c_{\nu,p}\) denote a \(\chi^2_\nu\) quantile. The event
\(c_{2n,\alpha/2}\le2\lambda W\le c_{2n,1-\alpha/2}\)
has probability \(1-\alpha\). Since \(W>0\) almost surely, dividing
and inserting the observed total \(w\) gives the exact rate interval
\[
  \left[\frac{c_{2n,\alpha/2}}{2w},\quad
        \frac{c_{2n,1-\alpha/2}}{2w}\right].
\]
For \(n=5\) and observed total wait \(w=15\) minutes,
\(c_{10,0.025}\approx3.247\) and
\(c_{10,0.975}\approx20.483\). The \(95\%\) interval is therefore
\[
  [0.1082,\,0.6828]\ \text{requests per minute}.
\]
Five observations leave considerable uncertainty around the point
estimate \(1/3\). The chi-square quantiles give unequal margins below
and above this estimate: approximately \(0.2251\) and \(0.3495\).
The exact coverage guarantee relies on independent exponential waits.

\section{Proportion intervals and test accuracy}
\label{sec:proportion-intervals}

For a fixed classifier evaluated on independently sampled examples,
let \(X_i=1\) when it predicts correctly and \(X_i=0\) otherwise.
The population accuracy is \(p\), and the observed accuracy is
\(\widehat p=k/n\), where \(k\) is the number correct.
The Bernoulli calculation gives standard error
\(\sqrt{p(1-p)/n}\). Substituting \(\widehat p\) gives the
familiar Gaussian approximation.

\begin{result}[Approximate Wald interval for a proportion]
For i.i.d. Bernoulli observations, the normal approximation gives
\[
  \widehat p\pm z_{1-\alpha/2}
     \sqrt{\frac{\widehat p(1-\widehat p)}n}.
\]
This approximation can be poor when successes or failures are rare.
It is not an exact interval, and its endpoints can fall outside \([0,1]\).
\end{result}

For \(k=84\) correct predictions out of \(n=100\),
\[
  \widehat p=0.84,\qquad
  \widehat{\operatorname{SE}}\approx0.0367,
  \qquad C_{\mathrm{Wald}}\approx[0.768,\,0.912].
\]
This quantifies uncertainty from sampling the test examples for the
already fitted classifier. It does not measure how much accuracy
changes when we retrain the model on a new training set.

\subsection{Why zero observed failures do not imply certainty}

If all \(20\) predictions are correct, then \(\widehat p=1\), and
the plug-in standard error is zero. The Wald formula returns \([1,1]\).
Twenty correct predictions provide evidence of high accuracy, but
cannot establish perfect population accuracy. Similarly, zero observed
successes do not establish \(p=0\). Clipping endpoints to \([0,1]\)
does not repair the underlying coverage problem.

A useful improvement keeps the candidate population probability
inside its own variance when inverting the approximate Gaussian event:
\[
  \frac{|\widehat p-p|}{\sqrt{p(1-p)/n}}\le z,
  \qquad z=z_{1-\alpha/2}.
\]
Squaring and rearranging gives
\[
  (n+z^2)p^2-(2n\widehat p+z^2)p+n\widehat p^2\le0.
\]
The two roots supply an interval.

\begin{result}[Wilson interval for a proportion]
Set \(z=z_{1-\alpha/2}\). The center and half-width are
\[
  a=\frac{\widehat p+z^2/(2n)}{1+z^2/n},\qquad
  h=\frac{z}{1+z^2/n}
       \sqrt{\frac{\widehat p(1-\widehat p)}n+\frac{z^2}{4n^2}}.
\]
The interval is \([a-h,a+h]\). It stays inside \([0,1]\) and usually
has much better coverage than the Wald interval near the boundaries.
Its \(1-\alpha\) confidence level is still approximate.
\end{result}

For \(84/100\), the \(95\%\) Wilson interval is approximately
\([0.756,0.899]\). For \(0/20\), it is \([0,0.161]\), and for
\(20/20\), it is \([0.839,1]\). In each boundary case the unobserved
outcome remains possible. Exact binomial intervals can provide a
guaranteed coverage level, generally with conservative coverage because
the count is discrete.

\begin{figure}[H]
  \centering\InfProportionCoveragePlot
  \caption{Actual coverage of nominal \(95\%\) Wald and Wilson
  procedures for \(n=20\). At each candidate \(p\), coverage is the
  exact binomial sum over counts whose interval contains \(p\).
  The jagged curves reflect the discrete count. Wilson improves the
  boundary behavior, but its coverage is not identically \(0.95\).}
  \label{fig:proportion-coverage}
\end{figure}

\section{Bootstrap and the empirical distribution}
\label{sec:bootstrap}

For a mean under a Gaussian model, we could derive a sampling law.
For a median, a quantile, or an F1 score, the sampling distribution
can be harder to calculate. We would like to repeat the experiment
many times, but usually have only one observed dataset. The bootstrap
uses this dataset to construct an approximation to the population,
then repeats sampling from that approximation.

\subsection{Resampling observations with replacement}

\begin{definition}[Empirical distribution and bootstrap sample]
Given observations \(x_1,\ldots,x_n\), the empirical distribution
\(\widehat F_n\) puts probability \(1/n\) on each recorded observation.
Repeated values contribute their combined probability.
A nonparametric bootstrap sample draws
\(X_1^*,\ldots,X_n^*\) independently from \(\widehat F_n\), conditional
on the observed dataset. Equivalently, draw indices
\[
  I_1,\ldots,I_n\ \text{independently and uniformly from }
  \{1,\ldots,n\},\qquad X_j^*=x_{I_j}.
\]
Sampling is with replacement, so an observation can appear several
times or be absent from a resample.
\end{definition}

For the waits \((1,2,3,4,5)\), the index sample \((1,1,2,4,5)\)
produces \((1,1,2,4,5)\), whose mean is \(2.6\). Another draw might
produce \((2,3,3,3,4)\), whose mean is \(3\). Both resamples have
the original sample size of five. Sampling all observations without
replacement would merely permute them and leave the mean unchanged.

For any statistic \(T\), recompute \(\widehat\theta^*=T(X^*)\)
on each resample. If the original estimate required fitting a model,
recomputing it requires fitting that model again. The resamples give
variation in the \emph{estimation rule}, rather than new independent
information about the population.

\subsection{Which distribution the bootstrap approximates}

An asterisk indicates randomness in resampling, conditional on the
recorded data. For a population parameter \(\theta\) estimated by
\(\widehat\theta\), the useful comparison is
\[
  \text{bootstrap errors }\widehat\theta^*-\widehat\theta
  \quad\text{versus sampling errors }\widehat\theta-\theta.
\]
The empirical population plays the role of the unknown population.
For a mean, its center is \(\bar x\), so bootstrap errors about
\(\bar x\) take the place of sampling errors about \(\mu\).
Similar reasoning applies to medians and other summaries when the
empirical population captures the features relevant to the statistic.
That approximation requires conditions and can fail, as
Section~\ref{sec:bootstrap-limits} explains.

For a mean we can check the calculation directly. Conditional on the data,
\[
  \Expect^*[X_j^*]=\bar x,\qquad
  \Var^*(X_j^*)=\frac1n\sum_i(x_i-\bar x)^2.
\]
Independence within a bootstrap sample therefore gives
\begin{equation}
  \Expect^*[\bar X_n^*]=\bar x,\qquad
  \Var^*(\bar X_n^*)=\frac1{n^2}\sum_i(x_i-\bar x)^2
  =\frac{n-1}{n}\frac{s^2}{n}.
  \label{eq:bootstrap-mean-variance}
\end{equation}
For \((1,2,3,4,5)\), \(\bar x=3\) and
\(\sum_i(x_i-\bar x)^2=10\). Thus the bootstrap mean has variance
\(0.4\) and standard error \(\sqrt{0.4}\approx0.632\).
The usual estimate \(s/\sqrt5=\sqrt{0.5}\approx0.707\) differs
by the factor \(\sqrt{(n-1)/n}\). The difference is appreciable at
\(n=5\) and tends to zero as \(n\) grows. A bootstrap approximation
does not guarantee an exact finite-sample sampling distribution.

\begin{figure}[H]
  \centering\InfBootstrapPlot
  \caption{Bootstrap means from one simulated dataset of \(40\)
  exponential waits. Its observed mean is \(4.495\) minutes. Each of
  \(20{,}000\) resamples contains \(40\) observations drawn with
  replacement from this same dataset. The dotted line marks the
  observed mean, and the dashed lines mark the \(2.5\%\) and
  \(97.5\%\) bootstrap quantiles. This is a conditional bootstrap
  distribution, rather than the true population sampling distribution.}
  \label{fig:bootstrap-errors}
\end{figure}

\subsection{Bootstrap standard error}

Choose a number of resamples \(B\). For each
\(b=1,\ldots,B\), draw \(n\) indices with replacement and recompute
\(\widehat\theta^{*(b)}\). Common choices such as \(B=10{,}000\)
give many replicates for estimating spread and central quantiles.
Record the random seed when reproducing a numerical illustration.

\begin{definition}[Bootstrap standard error]
Let \(\overline\theta^*=B^{-1}\sum_{b=1}^B\widehat\theta^{*(b)}\),
with \(B\ge2\). The bootstrap standard error estimate is
\[
  \widehat{\operatorname{SE}}_{\mathrm{boot}}
  =\sqrt{\frac1{B-1}\sum_{b=1}^B
    (\widehat\theta^{*(b)}-\overline\theta^*)^2}.
\]
It estimates the conditional standard deviation of the bootstrap
statistic, which in turn approximates the statistic's sampling
standard deviation when bootstrap assumptions hold.
\end{definition}

The divisor is \(B-1\) because we estimate spread across \(B\)
replicates about their replicate mean. The statistic inside each
replicate still uses a dataset of size \(n\). These are different
sample sizes with different roles. Increasing \(B\) reduces simulation
noise in the estimated bootstrap distribution. Increasing \(n\)
provides more information about the population.

\subsection{Percentile intervals and centered-error intervals}

\begin{definition}[Percentile bootstrap interval]
Let \(q_p^*\) denote the empirical \(p\)-quantile of the
\(B\) bootstrap estimates. The percentile interval is
\[
  [q_{\alpha/2}^*,\ q_{1-\alpha/2}^*].
\]
For a nominal \(95\%\) interval, retain the central \(95\%\) of
the bootstrap estimates. Its coverage is approximate.
\end{definition}

This interval can be asymmetric, following the shape of the resampled
statistic. It is easy to compute, but the fact that \(95\%\) of
bootstrap estimates lie between its endpoints does not by itself prove
\(95\%\) coverage of the population parameter. Bias and a poor
bootstrap approximation can distort coverage.

There is another construction that follows the centered-error
comparison explicitly. If \(a^*,b^*\) are the lower and upper
quantiles of \(\widehat\theta^*-\widehat\theta\), approximate the
sampling event by
\[
  a^*\le\widehat\theta-\theta\le b^*.
\]
Inverting gives \([\widehat\theta-b^*,\widehat\theta-a^*]\), or
\[
  [2\widehat\theta-q_{1-\alpha/2}^*,\quad
   2\widehat\theta-q_{\alpha/2}^*].
\]
This is the \emph{basic bootstrap interval}. It differs from the
percentile interval unless the relevant quantiles are symmetric
about \(\widehat\theta\). Neither construction is automatically
accurate for every statistic. A symmetric Gaussian approximation
using \(\widehat\theta\pm z\widehat{\operatorname{SE}}_{\mathrm{boot}}\)
is another option when its shape and bias assumptions are appropriate.

\Needspace{7\baselineskip}
For the simulated dataset in Figure~\ref{fig:bootstrap-errors},
\(\bar x\approx4.495\) minutes and \(s\approx3.747\) minutes.
The \(20{,}000\) bootstrap means give a standard error estimate
of \(0.586\) minutes and a percentile interval
\([3.399,5.682]\) minutes. We generated the original data from an
exponential population with mean five, so we can check that this
particular interval covers it. Checking one successful interval does
not establish the procedure's coverage. That would require repeating
the whole experiment with many fresh datasets and constructing a new
bootstrap interval for each one.

\section{Comparing two models on the same examples}
\label{sec:paired-bootstrap}

Suppose two fixed classifiers predict the same \(200\) independent
test examples. Their outcomes are
\begin{center}
\begin{tabular}{@{}lrr@{}}
\toprule
 & B correct & B wrong\\
\midrule
A correct & 150 & 30\\
A wrong   & 10  & 10\\
\bottomrule
\end{tabular}
\end{center}
Model A has accuracy \(180/200=0.90\), and B has accuracy
\(160/200=0.80\). The observed difference is \(0.10\).
Most examples are easy or difficult for both models, so treating the
two accuracy estimates as independent discards useful information.

Let \(A_i,B_i\) be the two correctness indicators on example \(i\),
and write \(D_i=A_i-B_i\). Then \(D_i=1\) on the \(30\)
examples only A gets right, \(D_i=-1\) on the \(10\) examples only
B gets right, and \(D_i=0\) on the other \(160\). The estimated
population accuracy difference is simply \(\bar D=0.10\).
The sample mean calculation applies to these \emph{paired differences}:
\[
  \widehat{\operatorname{SE}}(\bar D)
  =\sqrt{\frac1{n(n-1)}\sum_i(D_i-\bar D)^2}
  =\sqrt{\frac{0.20-0.10^2}{199}}\approx0.0309.
\]
The covariance formula in Chapter~\ref{chap:joint-distributions}
explains the same result:
\(\Var(A_i-B_i)=\Var(A_i)+\Var(B_i)-2\Cov(A_i,B_i)\).

\begin{result}[Paired bootstrap for a metric difference]
Keep each test example's label and both model predictions in one row.
For each replicate, sample \(n\) rows with replacement, using the
same indices for both models. Recompute both metrics and their
difference on the resampled rows. Apply the standard-error or
interval calculation to the resulting differences.
\end{result}

Resampling the two models independently would remove their within-example
dependence and estimate a different experiment. Pairing preserves the
fact that both models encounter the same resampled easy and difficult
examples.

With \(20{,}000\) paired resamples of this table, the bootstrap
standard error is approximately \(0.0310\) and the percentile
\(95\%\) interval is \([0.04,0.16]\). These values describe an
estimated accuracy improvement of \(4\) to \(16\) percentage points.

\begin{figure}[H]
  \centering\InfPairedPlot
  \caption{Paired bootstrap accuracy differences for the displayed
  \(200\)-example table. Each replicate draws the same test rows for
  both classifiers. The percentile endpoints describe uncertainty
  from sampling test examples for these fixed fitted models.}
  \label{fig:paired-bootstrap}
\end{figure}

For precision, recall, and F1, retain the label and predictions in each
row, then recompute the confusion counts within each resample. For
example, precision is \(\mathrm{TP}/(\mathrm{TP}+\mathrm{FP})\),
recall is \(\mathrm{TP}/(\mathrm{TP}+\mathrm{FN})\), and
\[
  \mathrm{F1}=\frac{2\mathrm{TP}}{2\mathrm{TP}+\mathrm{FP}+\mathrm{FN}}.
\]
These are ratios of counts, so averaging hypothetical per-example
F1 scores would not reproduce the metric. With rare positive examples,
some resamples may have a zero denominator. Its handling must be
specified and the resulting interval may be unstable.

An interval for an accuracy difference that lies above zero is evidence
of an improvement for the stated test population and sampling procedure.
It does not imply an improvement for every example. If the interval
includes zero, the observed experiment may be too imprecise to resolve
the difference. Selecting the two models by repeatedly comparing many
candidates on this same test set changes the uncertainty question.

\section{What resampling can and cannot recover}
\label{sec:bootstrap-limits}

\subsection{Small samples and missing outcomes}

The empirical distribution contains only observed values. A very small
dataset may omit an important part of the population, especially a
rare but influential outcome. No number of ordinary resamples can
recover an event that has zero probability under the empirical
distribution. For example, bootstrapping twenty observed failures
as Bernoulli outcomes gives failure rate one in every resample.
The resulting zero spread is the same warning we encountered with
the Wald interval.

Extreme-value statistics provide a sharper counterexample. For
\(X_i\sim\operatorname{Uniform}(0,\theta)\), let
\(M=\max_iX_i\), the endpoint MLE from
Section~\ref{sec:likelihood-estimation}. Every ordinary bootstrap
maximum satisfies \(M^*\le M\). But the true endpoint satisfies
\(\theta>M\) with probability one. Thus the upper endpoint of a
percentile interval made from \(M^*\) is at most \(M\), and the
interval covers \(\theta\) with probability zero. A sampling model
can still support inference about \(\theta\), but ordinary empirical
resampling fails for this endpoint problem.

\subsection{The independent unit determines what to resample}

If several measurements come from one user, independently resampling
measurements loses their shared user effect. When users are independent
sampling units, a cluster bootstrap resamples users and keeps their
measurements together. For time series, resampling suitable consecutive
blocks can retain some temporal dependence. Block length and assumptions
about the process matter. Ordinary i.i.d. row resampling does not
automatically apply to these settings.

The sample must also represent the target population. Repeating a
biased sampling scheme or bootstrapping a nonrepresentative dataset
does not remove its selection bias. The same problem arises when the
deployment population differs from the test population. Heavy tails
can also make bootstrap approximation unreliable, particularly for
means with infinite variance.

\subsection{Parameter uncertainty and training uncertainty}

For an interval about a fitted parameter, resample the independent
training units and fit the estimator again on each resample. For
regression, preserve each feature vector and its response as a row.
This approximates randomness in drawing those training rows under
the stated sampling model. A design that holds feature values fixed
requires a different resampling construction, such as resampling
modeled errors, with its own assumptions.

The paired test bootstrap in Section~\ref{sec:paired-bootstrap}
holds the trained models fixed. It describes test-sample uncertainty.
To include variation from training data, random initialization, or
hyperparameter selection, the repeated experiment must also repeat
the corresponding training or selection steps. Recomputing only
test metrics cannot estimate those additional sources of variation.

\endgroup
